\section{A topological correspondence for every flag}
\label{top:section}

Contour convergence initially relates exploration times.  We need a
correspondence on the whole discrete surface which preserves separation
and intersections.  This section constructs it without estimating the
length of any connecting walk.  Throughout the finite-volume argument we
work on a coupling on which the contours converge uniformly, write
$Z_n=(L_n,R_n)$ for the rescaled contours, and write
$\eta:[0,1]\to S$ for the limiting exploration.  Equip the reference
topological sphere $S$ with a fixed spherical metric $\rho$ in any of its
conformal coordinates.  The choice of this coordinate is immaterial to
the deterministic assertions below.

\subsection{Scalar trees and the continuous projection}

For a continuous real function $c$ on a compact time interval, put
\begin{equation}\label{top:tree-pseudometric}
 d_c(s,t)=c(s)+c(t)-2\min_{[s\wedge t,s\vee t]}c.
\end{equation}
The equivalence classes for $d_c=0$ are the vertices of its contour tree.
We use the exact primal-tree and dual-tree incidence description in
Lemma~\ref{enc:incidence}.  Heights of a discrete tree edge in this
section have the same rescaled units as $L_n,R_n$; its graph length is
still one whenever graph costs are subsequently considered.

\begin{lemma}[Continuity of the scalar trees]\label{top:scalar-trees}
If $c_n\to c$ uniformly, then
\[
 \sup_{s,t}|d_{c_n}(s,t)-d_c(s,t)|
 \le 4\|c_n-c\|_\infty.
\]
Suppose $s_n\to s$, $t_n\to t$, and $d_c(s,t)=0$.
Every contour representative of every vertex on the discrete
$c_n$-tree path between the vertices at $s_n,t_n$ has distance
$d_{c_n}$ tending to zero from the vertex at $s_n$, uniformly over those
representatives.  Consequently, if $c$ is $L$ or $R$, their images under
$\eta$ tend uniformly to $\eta(s)$.

For $c_n=R_n$, take the faces represented by the vertices of this
dual-tree path.  The union of their primal boundary walks is a connected
primal subgraph.  All its corner representatives have $\eta$-images
tending uniformly to $\eta(s)$.
\end{lemma}

\begin{proof}
The displayed inequality follows by bounding the two endpoint errors and
the error of the minimum.  A vertex $u$ on a tree geodesic from $a$ to $b$
satisfies $d(a,u)\le d(a,b)$.  The exact contour-tree identity therefore
bounds all the asserted distances by
$d_{c_n}(s_n,t_n)$, up to one rescaled edge if a cut lies inside a step.
This tends to zero.  If the assertion about $\eta$ were false, choose a
representative $u_n$ violating it and then a convergent subsequence of
$u_n$.  Its limit $u$ satisfies $d_c(s,u)=0$, so
$\eta(u)=\eta(s)$ by the scalar chord rule, contradicting continuity of
$\eta$.

Consecutive faces of a dual-tree path share the primal edge crossed by
that tree edge.  Their entire boundary walks therefore have connected
union, including when a boundary edge or vertex is visited repeatedly.
A corner of any of these faces represents its dual-tree vertex, so the
same distance bound applies to all its representatives.  This proves
the last assertion.  No bound on a face degree is involved.
\end{proof}

The tour supplies finitely many time representatives for each flag
incidence, as specified in Lemma~\ref{enc:incidence}.  We may use
either endpoint of an exploration step; changing that choice changes the
time by at most one mesh interval.  A representative of a primal vertex
means any tour occurrence of that vertex, and similarly for a dual
vertex.  The occurrences at a loop or bridge remain distinct.

\begin{proposition}[Projection of the complete flag surface]
\label{top:projection}
There are continuous maps $P_n:X_n\to S$ and numbers
$\varepsilon_n\to0$ such that the following hold simultaneously.
\begin{enumerate}[label=(\roman*)]
\item For every flag $T$, every $x\in T$, and every tour representative
associated to any of the corners of its quadrangle,
\[
 \rho(P_n(x),\eta(t))\le\varepsilon_n.
\]
The same conclusion holds for any other occurrence of an incident primal
or dual vertex.
\item The $P_n$-images of individual flags, original edges, original
closed faces, and complete primal or dual incidence stars have maximum
diameter tending to zero.
\item The sets $P_n(V(M_n))$ become dense in $S$.
\end{enumerate}
These conclusions hold for the original flag complex, with all repeated
incidence occurrences retained.
\end{proposition}

\begin{proof}
First choose one representative for every primal vertex and every dual
vertex.  Representatives of a given primal vertex have zero
$L_n$-tree distance; representatives of a given dual vertex have zero
$R_n$-tree distance.  Uniform contour convergence, compactness of the
time interval, and the chord rule show that the maximum disagreement
between their $\eta$-images tends to zero.  This is a uniform statement:
a failing sequence of vertices and two of their occurrences would have
convergent time representatives and would contradict the limiting chord
rule.

At a tour cut the primal and dual corner states have that same time
representative.  In one quadrangle its four corner states are linked
successively by equality of a primal or a dual vertex, together with the
two consecutive-cut errors of its exploration triangles.  The preceding
uniform assertion and continuity of $\eta$ show that all these
representatives agree up to a uniform error tending to zero.  It follows
that the chosen representatives of the primal vertex and face center
of every flag agree to this accuracy.  Assign its edge midpoint the
image of one fixed corner occurrence of that edge.  The preceding
quadrangle argument gives the same conclusion for this assignment.

Identify $S$ with the unit sphere in $\R^3$.  On every flag send a point
with barycentric coordinates $(a_1,a_2,a_3)$ to
\[
 \frac{a_1 y_1+a_2 y_2+a_3 y_3}
 {|a_1 y_1+a_2 y_2+a_3 y_3|},
\]
where $y_1,y_2,y_3$ are the assigned images of its vertices.  For all
large $n$ the three images lie in a common ball of radius less than
$\pi/4$, so the denominator is nonzero.  These formulas agree on glued
sides and give a continuous map on the flag quotient.  Their images
stay in the same small spherical ball.  This proves (i) and the flag
assertion in (ii).  Every flag in a primal or dual star shares the
image of its central vertex, so its entire star has diameter at most
twice the maximum flag diameter.  The assertions for edges and faces
follow as well.  Finally the tour meshes become dense in $[0,1]$,
$\eta$ is onto, and every tour position has a primal corner vertex
with the stated small error.  This proves (iii).
\end{proof}

\begin{lemma}[Protection by complete time preimages]
\label{top:protection}
Let $K\subset U\subset S$, where $K$ is compact and $U$ is open.
There are a neighborhood $K'$ of $K$ with $\overline{K'}\subset U$ and a
finite union $I$ of open time intervals such that
\begin{equation}\label{top:time-protection}
 \eta^{-1}(\overline{K'})\subset I,
 \qquad \overline I\subset\eta^{-1}(U).
\end{equation}
Intervals are interpreted on the time circle when necessary.  If
$p_*\notin\overline U$, they are strict subintervals of $(0,1)$.
For all large $n$, every flag whose projection meets $K$, and the full
incidence stars of its three vertices, have all their tour
representatives in $I$.

In the bilateral setting the same assertion holds for bounded compact
spatial protections, with finitely many bounded time intervals.
\end{lemma}

\begin{proof}
Choose two successive compact neighborhoods of $K$ in $U$.
The inverse image of the smaller closed neighborhood is compact and
lies inside the open inverse image of the larger one.  Cover it by
finitely many intervals whose closures lie in that open inverse image.
This proves \eqref{top:time-protection}.  If a sequence of omitted tour
representatives belonged to stars of flags meeting $K$,
Proposition~\ref{top:projection} would force their $\eta$-images into
an arbitrarily small neighborhood of $K$.  They would eventually lie
in $\eta^{-1}(\overline{K'})\subset I$, a contradiction.  Avoiding
$p_*$ keeps the inverse images away from $0,1$.
In the plane, transience of $\eta$ makes the inverse image of every
compact set compact; the same finite-cover proof then applies.
\end{proof}

\subsection{Joining arbitrary lifts with spatial clearance}

\begin{lemma}[Local joining]\label{top:local-joining}
For every $\epsilon>0$ there are $\delta>0$ and $n_0$ such that, for
$n\ge n_0$, $x\in S$, and primal vertices $u,v$ satisfying
\[
 P_n(u),P_n(v)\in B_\rho(x,\delta),
\]
there is a primal graph walk from $u$ to $v$ whose entire projection is
contained in $B_\rho(x,\epsilon)$.
The analogous assertion holds for arbitrary prescribed points of
$X_n$, using continuous paths.  The choice is uniform over the
prescribed lifts.  Its conclusion makes no assertion about the number
of edges in the walk.
\end{lemma}

\begin{proof}
First suppose that representatives $s_n,t_n$ of $u_n,v_n$ converge to
$s,t$ with $s\sim_L t$.  The primal-tree path joins these vertices and
has image tending to the single point $\eta(s)$ by
Lemma~\ref{top:scalar-trees} and Proposition~\ref{top:projection}.
If instead $s\sim_R t$, join the incident dual vertices by their
dual-tree path.  The union of the primal boundaries of its faces is
connected, contains both primal corner vertices, and has the same
vanishing image diameter.  A walk in this connected subgraph supplies
the required primal connection.

For arbitrary $s,t$ with $\eta(s)=\eta(t)$, choose a finite chain of
$L$- and $R$-chords joining them, using
Theorem~\ref{enc:peanosphere}.  For each intermediate time choose its
nearest discrete cut and its primal corner vertex.  The preceding
construction for each link gives finitely many walks whose union tends
to $\eta(s)$; concatenate them.

If the uniform assertion failed, there would be $n_j\to\infty$,
$\delta_j\downarrow0$, points $x_j$, and two prescribed vertices
with images in $B(x_j,\delta_j)$ admitting no walk in
$B(x_j,\epsilon)$.  Choose time representatives and pass to convergent
subsequences $x_j\to x$, $s_j\to s$, $t_j\to t$.
Proposition~\ref{top:projection} gives
$\eta(s)=\eta(t)=x$.  The finite-chain construction just proved
produces the forbidden walks.  For an arbitrary point of a flag, first
join it inside that flag to its primal vertex.  The extra pieces have
vanishing projected diameter by Proposition~\ref{top:projection}.
\end{proof}

To specify an ordered following test for a path $c:[0,1]\to S$,
choose $0=t_0<\cdots<t_r=1$, open ports $U_j$ containing $c(t_j)$,
and open tubes $V_j$ containing $c([t_{j-1},t_j])$.
A follower visits these ports in order, with its intervening pieces
in the corresponding tubes.  The compact reference pieces and the
chosen smaller port neighborhoods have positive clearance inside
these open sets.  A port specifies a visit; a transverse crossing
will require the additional side and clearance conditions below.

\begin{corollary}[Following a finite collection of paths]
\label{top:path-lifting}
Let $c_1,\ldots,c_m$ be continuous paths on $S$ and let prescribed
lifts of their endpoints have the corresponding projected limits.
For every finite ordered following test as above, all sufficiently
large $n$ admit lifted paths whose projections satisfy that test.
Primal endpoint lifts admit
primal walks.  Endpoints which coincide in the reference collection
may be assigned one common prescribed lift.
\end{corollary}

\begin{proof}
Apply Lemma~\ref{top:local-joining} with a spatial error smaller than
the clearance of the chosen finite requirements.  Uniform continuity
gives a finite time partition for each $c_i$ on whose pieces its
oscillation is smaller than the resulting $\delta/3$.  Choose primal
representatives near the partition points using
Proposition~\ref{top:projection}(iii), retaining the prescribed endpoints.
Join successive representatives by the lemma.  Subdivide at the
intermediate ports before making these choices.  Concatenation proves
the ordered assertion, and the same construction works for all the
finitely many paths at once.
\end{proof}

\subsection{From local joining to sphere homeomorphisms}

Local joining initially supplies connected branch sets, rather than
disjoint arcs meeting only at their endpoints.  We use a cubic
skeleton because a minimal tree joining three attachments has just
one trivalent junction after its outgoing paths are attached.
Pruning therefore turns the branch sets into a subdivision, to which
Whitney's embedding theorem and disk extension can be applied.

\begin{lemma}[Lifting a cubic skeleton]\label{top:skeleton-lift}
Let $G\subset S$ be an embedded finite simple cubic graph with
piecewise smooth edges.  For every $\epsilon>0$, all sufficiently large
$n$ contain a primal subgraph $G_n$ which is a subdivision of $G$, and
there is a graph homeomorphism $h_n:G_n\to G$ such that
\[
 \sup_{x\in G_n}\rho(P_n(x),h_n(x))<\epsilon.
\]
If $G$ is $3$-connected, its facial cycles and those of $G_n$ correspond
under this homeomorphism, up to one common reversal of orientation.
\end{lemma}

\begin{proof}
Take pairwise disjoint small closed junction disks about the vertices
of $G$.  In each disk its three incident edges are three disjoint stubs
except for their common endpoint.  Choose three further intermediate
points on the stubs, and narrow disjoint ribbons along the remaining
edge portions.  The ribbons of different edges are disjoint outside
the junction disks.  Inside a junction disk, choose the three ribbon
ends disjoint, so each meets only its own stub neighborhood.  Make
all these neighborhoods small enough for the desired error.

Assign one primal lift to every graph vertex.  By
Corollary~\ref{top:path-lifting}, lift its three initial stubs from
that common lift to the three intermediate points.  Their union is a
connected branch set $B_v$ within the chosen junction neighborhood.
The branch sets belonging to distinct vertices are disjoint, since
their projected neighborhoods are disjoint.  Lift each remaining
edge portion between the corresponding stub endpoints.  The error
may be chosen so small that different connecting walks can intersect
only within branch sets at a common endpoint, and no connecting walk
meets any other branch set.  In particular, near a junction each
connector meets only the end of its own stub; all such contacts occur
before the prescribed marks delimiting the middle of that connector.

For an edge $vw$, take the portion from its last visit to $B_v$
before its first subsequent visit to $B_w$.  Erase loops from that
finite walk.  The resulting simple paths have interiors disjoint
from every branch set and from each other.  At a branch set retain
a minimal finite tree joining its three attachment points.  After
the three outgoing paths are attached, this tree has exactly one
vertex of degree three and every other vertex of degree two; this
also covers coincident attachment points and the case where one
attachment lies between the other two.  Removing unused branches
therefore gives an actual subdivision of $G$.  The use of three
attachments is the reason for choosing a cubic skeleton.

Here is the order control needed to define $h_n$.  Give each edge
ribbon a continuous longitudinal coordinate, equal to its edge
parameter on the central arc.  In the path-lifting construction use
successive coordinate intervals of length less than $\sigma$, and
choose the spatial errors so that each joining piece stays within
$\sigma$ of its prescribed interval.  If a loop-erasure removes a
portion between two occurrences of one vertex, the two original
parameters differ by at most $3\sigma$: their common projected
point has a single longitudinal coordinate within the required
error of both.  Truncation at the branch sets changes only the
initial or final stub.  Assign increasing central-arc parameters to
the surviving vertices, and interpolate strictly increasingly
along the surviving edges.  The skipped parameter intervals have
oscillation tending to zero as $\sigma\downarrow0$.
Within a junction neighborhood map the three remaining tree arms
to the three initial edge pieces.  This gives a graph homeomorphism
whose error is bounded by the ribbon widths, junction diameters,
and those vanishing oscillations.  Choose them in that order to
make the total error smaller than $\epsilon$.

Finally a subdivision induces the same embedding of the suppressed
graph.  Whitney's theorem says that two embeddings of a finite
simple $3$-connected planar graph in the sphere are equivalent up
to a sphere homeomorphism \cite{Whitney1932}.  Its hypotheses hold
for $G$, and hence the identified facial cycles agree, with at most
one global orientation reversal.
\end{proof}

\begin{theorem}[Homeomorphisms with vanishing projection error]
\label{top:homeomorphism}
There are sphere homeomorphisms $H_n:X_n\to S$ such that
\begin{equation}\label{top:homeomorphism-error}
 \sup_{x\in X_n}\rho(H_n(x),P_n(x))\longrightarrow0.
\end{equation}
In particular,
\begin{equation}\label{top:homeomorphism-mesh}
 \max_{T}\diam_\rho H_n(T)\longrightarrow0,
\end{equation}
and the images under $H_n$ of every flag and all its tour
representatives have vanishing maximum disagreement.  The maps may
preserve or reverse orientation, with a common sign on each sphere.
They may be chosen measurably in the contour coupling.  The assertion
continues to hold when that coupling retains any additional random
observations.
\end{theorem}

\begin{proof}
Fix a small $\epsilon>0$.  Triangulate the reference sphere with
simplicial triangles of diameter less than $\epsilon/10$, and take
its embedded dual graph $G$.  Draw dual edges through adjacent
triangle interiors.  Its faces lie in stars of the original vertices,
so their diameters are less than $\epsilon$.
This graph is cubic and $3$-connected.  For the latter assertion,
delete any two dual vertices, or equivalently the corresponding
two closed triangles.  Their union has connected complement in
the sphere, whether the triangles are disjoint, meet at a vertex,
or share an edge.  A path in that complement between two remaining
triangle interiors can be perturbed away from triangulation
vertices; the triangles it successively crosses give a dual path
avoiding the two deleted vertices.  Thus deleting fewer than three
dual vertices does not disconnect the dual graph.  The simplicial
condition also ensures that the dual has neither loops nor multiple
edges.

Apply Lemma~\ref{top:skeleton-lift} with an error much smaller than
$\epsilon$.  Denote by $F_n$ the face of $G_n$ corresponding to a
face $F$ of $G$.  Choose $z_F$ so that
$\overline F\subset B(z_F,\epsilon)$.  We claim that, for all
large $n$,
\begin{equation}\label{top:small-faces}
 P_n(\overline{F_n})\subset B(z_F,3\epsilon).
\end{equation}
The boundaries already satisfy this assertion by the graph
homeomorphism estimate.  If the claim fails in the interiors,
pass to a sequence $x_n\in F_n$ with
$P_n(x_n)\to y\notin B(z_F,3\epsilon)$.  Choose a vertex $v$ of
$G$ outside $B(z_F,3\epsilon)$; such a vertex exists for sufficiently
small $\epsilon$, because the dual graph has small faces and covers
the sphere.  The complement of
$\overline B(z_F,2\epsilon)$ is connected, so join $y$ to $v$ by
a path with clearance from $B(z_F,3\epsilon/2)$.  Lift this path
from the prescribed point $x_n$ to the junction lift of $v$, using
Corollary~\ref{top:path-lifting}.  Its projection avoids
$B(z_F,3\epsilon/2)$ for all large $n$, whereas the projection of
$\partial F_n$ is contained in that ball.  The lifted path
therefore avoids $\partial F_n$.  Its last point is outside
$\overline{F_n}$, since $v$ is not incident to the corresponding
face.  This contradicts the Jordan separation of $\partial F_n$.
There are finitely many faces, so \eqref{top:small-faces} holds
for all of them simultaneously.

Each face of a $3$-connected sphere embedding is a Jordan disk.
The graph homeomorphism $h_n$ therefore extends across every face
to a homeomorphism $H_{n,\epsilon}:X_n\to S$.
For completeness, a homeomorphism between the boundary circles of
two closed disks extends by identifying each disk with the standard
closed disk and extending the resulting circle homeomorphism
radially.  The Jordan--Sch\"onflies theorem gives these disk
identifications; here all boundaries are finite polygonal curves
in triangulated surfaces.  The facial correspondence makes the
extensions agree on common edges.  For $x\in\overline{F_n}$,
both $H_{n,\epsilon}(x)\in\overline F$ and
$P_n(x)\in B(z_F,3\epsilon)$, and hence
\[
 \sup_{x\in X_n}\rho(H_{n,\epsilon}(x),P_n(x))
 \le 4\epsilon.
\]

Take $\epsilon_k\downarrow0$.  At each $k$ the construction works
eventually; increase its threshold $n_k$ so that the thresholds are
strictly increasing, and use $H_{n,\epsilon_k}$ for
$n_k\le n<n_{k+1}$.  This proves
\eqref{top:homeomorphism-error}.  Equation~\eqref{top:homeomorphism-mesh}
follows from Proposition~\ref{top:projection} and the triangle
inequality.

All finite walk choices can be made first in a fixed enumeration.
The graph and its disk extensions can likewise be chosen among
piecewise linear maps after finite subdivisions, with rational
coordinates in fixed spherical charts.  These form a countable
collection; the polygonal disk construction and an arbitrarily
small perturbation give a member satisfying any strictly weakened
error bound above.  Choose the first such member.  Its uniform
error is measurable by taking a supremum over a countable dense
set in each flag.  Choosing the largest available accuracy among
the first $n$ accuracies gives a measurable version of the
diagonal construction.  Finally, the whole argument is
deterministic on the event of contour convergence and the
peanosphere assertions, so retaining other observations does not
alter it.
\end{proof}

\subsection{Protected bilateral windows}

The analogous plane statement is local.  In particular, it does not
require a conformal uniformization of an infinite discrete map, or the
assertion that an arbitrarily capped limiting contour has a sphere
quotient.

\begin{proposition}[Closing a protected window]
\label{top:protected-window}
Work in a bilateral contour coupling with locally uniform convergence
to its peanosphere contour.  Given finitely many compact spatial
tests and open buffers about them, one may choose a bounded time
window, and then a finite spherical tree/cotree sewing, with the
following properties for all large $n$.
\begin{enumerate}[label=(\roman*)]
\item Every flag star needed by the tests is unchanged, including its
cyclic order and all original primal incidences.
\item On a neighborhood of these flags there are continuous projections
with the conclusions of Proposition~\ref{top:projection} and
Lemma~\ref{top:protection}, and the local joining conclusion of
Lemma~\ref{top:local-joining} holds with arbitrary fixed clearance.
\item The finite sewing admits homeomorphisms into the reference
one-point compactification for which the projection error tends to
zero over the tests.  The inverse images of the tested compact sets
are contained in the unchanged part.  Consequently both directions
of every strict crossing or separation test take place in the
original discrete surface.
\end{enumerate}
The construction can be made simultaneously for a countable exhaustion
of compact tests and their finite refinements, by diagonalization.
\end{proposition}

\begin{proof}
Transience makes the full time preimage of a spatial compact bounded.
Choose a large disk containing the tests and all their buffers, then a
larger disk for the finite joining routes below.  Both Brownian
coordinates are simultaneously below any prescribed pair of levels at
arbitrarily distant times in either direction.  Indeed, after an
invertible linear change of coordinates this follows from planar
Brownian recurrence applied to a disk contained in the corresponding
open cone.  Choose cuts $a<b$ beyond the complete time preimages in
use, at which both coordinates are strictly below all their heights.
Rounding the cuts to lattice times preserves these strict inequalities
for all large $n$.

No scalar chord from a protected time can cross either cut: it would
have to pass below its own endpoint height.  Thus every protected
primal vertex, dual vertex, matched quadrangle, and their full stars
are determined inside $(a,b)$.  Translate the resolved axial walk on
$[a,b]$ into the strictly positive quadrant, precede it by a monotone
axial path from the origin to its initial point, and follow it by a
monotone axial path from its terminal point to the origin.  This is a
finite nonnegative resolved excursion and hence a finite spherical
tree/cotree sewing.  The appended steps and all incidence classes
which they affect are below the protected heights.  The exact
tree/cotree incidence rule therefore leaves the protected full stars
unchanged.  This proves (i).

The projection and scalar-path proofs apply in any compact time
interval containing the complete preimages involved.  Strict lower
cuts confine all the relevant discrete representatives to that
interval.  A finite chain of limiting chords uses only finitely many
times, and its discrete tree paths can be read in a bounded enlargement.
The sequential compactness proof of local joining consequently gives
(ii), uniformly on each compact spatial buffer.  There is no need to
bound the length of one of these paths.

We give the exterior argument for (iii), since it prevents the cap
from entering a tested region.  Choose a finite simple cubic
$3$-connected planar skeleton $G$ whose bounded faces have diameter
less than $\epsilon$ throughout the large disk, and whose outer face
misses that disk.  Such a graph is the dual of a fine triangulation
of a disk completed on the sphere by a vertex at infinity; the dual
edges adjacent to that vertex can be drawn in a collar of the disk.
The proof of $3$-connectivity in Theorem~\ref{top:homeomorphism}
applies without change.  Lift its stubs and edge ribbons using (ii).
All these finite walks and their complete incidence stars are in a
bounded time protection.

Take $T$ larger than every time in this protection and so large that
both limiting tour tails $(-\infty,-T]\cup[T,\infty)$ avoid a disk
containing $G$ and its ribbons.  Retain two tour positions near
$-T,T$ as exterior sentinels.  For each bounded face of $G$, join each
sentinel to a graph vertex not incident to that face by a spatial
path avoiding a slightly enlarged disk about that face.  Retain also
a path between the two sentinels outside the disk containing $G$.
There are finitely many such routes.  Enlarge the time protection
to contain their lifted routes before choosing the final lower
cuts $a<-T<T<b$.

The tour strip from $T$ to $b$ is connected to the appended terminal
piece, and the strip from $a$ to $-T$ is connected to the appended
initial piece.  These strips avoid the lifted skeleton: an
intersection would be an incidence in one of its full stars, whose
representatives have already been retained between $-T$ and $T$.
All cap-affected flags are connected to these strips through the
appended tour, including any old exploration triangle whose partner
is supplied by an appended step.  Thus the cap and both strips lie
in one face of the lifted skeleton.  The sentinel routes show that
this face is not any of the bounded corresponding faces: a route
to a vertex not incident to such a face avoids its boundary, by its
projection clearance.  Whitney's facial correspondence therefore
identifies it as the outer face.  Equivalently, the cap changes
only that exterior component; every other complementary disk and
all its incidence data are the original ones.

For a bounded face, the connected-exterior proof of
\eqref{top:small-faces} now applies.  Its points have time
representatives within the fixed window, so a violating sequence
has a convergent representative.  The required path away from its
small boundary can be lifted in the original map by (ii); if that
path reaches a cut-affected region, it has already reached the outer
component, which gives the same separation contradiction.  Hence
the entire projected face lies in a $3\epsilon$ disk.  Extend the
graph homeomorphism over all bounded faces and extend arbitrarily
over the outer face.  Since the tested compact sets avoid the
outer face and its collar, their inverse images are in the unchanged
part, with error at most $4\epsilon$.  Let $\epsilon\downarrow0$.

All windows, skeletons, and routes just used are finite at each
stage.  Enumerate the desired compact tests and accuracies, retain
the first $k$ at stage $k$, and enlarge their windows before choosing
the cuts.  Choosing the discrete index after that finite stage
gives the asserted simultaneous diagonal version.
\end{proof}

\subsection{Strict intersections and primal separators}

\begin{corollary}[Transverse crossings]\label{top:transverse}
Fix a Jordan quadrilateral in a reference chart.  Prescribe one
crossing between its first and third sides and another between its
second and fourth sides, with endpoint ports beyond those sides,
and with positive clearance from the inappropriate sides.  Discrete
surface paths satisfying sufficiently accurate versions of these
requirements intersect.  If both are primal graph walks, they
have a common primal vertex.  The assertion holds simultaneously
for every fixed finite collection of such strict tests, on the
sphere or in a protected bilateral window.
\end{corollary}

\begin{proof}
Use $H_n$ to transfer the two walks to the reference sphere.  Its
uniform error is smaller than all the fixed port clearances for
large $n$.  Each walk then has a subpath crossing the indicated
opposite sides of a common quadrilateral.  Two such transverse
paths intersect: a simple subarc of the first joining its two
sides separates the other two sides, by the Jordan theorem.
Their inverse images intersect as well.  The interiors of
distinct edges in an embedded primal graph are disjoint, so an
intersection of primal walks gives a common vertex, also when
one of their edges is a loop or they share an edge.  Apply
Proposition~\ref{top:protected-window} in the bilateral case.
\end{proof}

\begin{corollary}[Primal surrounding circuits]\label{top:circuits}
Let $B_0\Subset B_1\Subset B_2\Subset B_3$ be concentric closed
disks in a planar reference chart.  Suppose a closed primal walk
has projection in $\operatorname{int}(B_2)\setminus B_1$ and
nonzero winding about $B_0$.  For all sufficiently accurate
correspondences it contains a simple primal circuit $C_n$ with
the following properties.  The closed side $\Omega_n$ of $C_n$
opposite the exterior of $B_3$ contains every lift whose projection
lies in $B_0$, is disjoint from every lift projecting outside $B_3$,
and satisfies
\[
 P_n(\Omega_n)\subset B_3.
\]
The circuit uses a subset of the walk's edges with multiplicities
discarded; for nonnegative edge costs its cost is at most the
cost of that walk.  The conclusion holds in a protected bilateral
window as well.
\end{corollary}

\begin{proof}
For large $n$, the $H_n$-image of the walk remains in a slightly
enlarged annulus between $B_0$ and $B_3$.  The short pointwise
homotopy from $P_n$ to $H_n$ avoids $B_0$ and an exterior point,
so the winding number is unchanged.  Decompose the finite closed
walk at repeated vertices into simple circuits, deleting immediate
backtracks.  Winding numbers add in this decomposition.  At least
one simple circuit has nonzero winding about $B_0$; its embedded
image is a Jordan curve, and therefore it separates all of $B_0$
from the exterior of $B_3$.  Its bounded side lies inside $B_3$,
because the connected exterior of $B_3$ avoids the circuit.
Keeping a smaller fixed error than the clearances between the
four disks gives precisely the assertions about all $P_n$-lifts
and about $P_n(\Omega_n)$.  The edge and cost assertion follows
from the circuit decomposition.  The protected-window
homeomorphism proves the local version.
\end{proof}

\begin{lemma}[Finite separator meshes from open ports]
\label{top:separator-mesh}
For every $\epsilon>0$ and every finite set of excluded reference
points, there is a finite list of reference paths avoiding those
points, with strict ordered ports and tubes, with this property:
any primal walks satisfying all the prescribed following tests
contain a connected primal subgraph whose complementary components
have $H_n$-image diameter at most $\epsilon$.
The subgraph uses only the supplied walks.  In particular its
total edge cost is at most the sum of their costs.

If a geodesic metric $D$ inducing the topology of $S$ is specified,
the reference paths can all be chosen to have finite $D$-length.
There is no cost requirement on joins from prescribed discrete
endpoints in this statement.
\end{lemma}

\begin{proof}
Choose a cubic $3$-connected skeleton with face diameter much
smaller than $\epsilon$, perturbing it to avoid the excluded
points.  Put disjoint small junction disks at its vertices and
narrow disjoint edge ribbons between them.  In each junction
disk prescribe four paths following the sides of a small square,
extended past its corners so that adjacent paths have strict
transverse crossings there.  Prescribe edge paths which run
from the central hole of one junction square through its edge
ribbon to the central hole of the other.

Corollary~\ref{top:transverse} stitches the four square paths
into a closed walk of nonzero winding about its central hole.
To see the winding directly, cut consecutive paths at their
forced intersections in the disjoint corner boxes and retain
the intervening side portions.  Their concatenation follows
the square once, and a homotopy within its four thin side
corridors preserves that winding.  Corollary~\ref{top:circuits}
then supplies a simple junction circuit.  Every incident edge
walk starts inside its central hole and exits its outer disk,
so it meets that circuit.  Thus all paths at a junction belong
to one connected branch set, without specifying their lattice
endpoints in advance.

Truncate each edge walk between its last visit to the first
branch set and its first visit to the second, and erase loops.
Different edge corridors are disjoint away from their assigned
junctions.  As in Lemma~\ref{top:skeleton-lift}, minimal trees
joining the three attachments produce a subdivision of the
cubic skeleton.  Its edges follow the reference edges in
order.  Under the actual homeomorphism $H_n$ it is an embedded
graph on the reference sphere with arbitrarily small tracking
error.  Whitney's facial correspondence and the connected
exterior of each small face, exactly as in the proof of
\eqref{top:small-faces}, imply that every complementary
component has diameter at most $\epsilon$.  Keeping the entire
connected union of the supplied walks only subdivides these
components.  No additional edges have been used, which proves
the cost assertion.

Finally a continuous path in an open set can be followed through
any fixed finite ordered ports by a finite-length $D$-path.
To verify this, cover its compact image by finitely many
$D$-balls whose doubled balls lie in the open set, subdivide
the parameter interval so successive selected points lie in
one such ball, and join them by short $D$-geodesics.  Those
geodesics remain in the doubled balls.  Applying this to the
finite path list, with smaller tubes and the same strict
ports, gives finite $D$-length reference paths with the
identical topological requirements.
\end{proof}
